%% RNAAS manuscript for UHI v5.2-final -- revision r3 (r2 + Christiansen et al. 2025 archive citation, 2026-11-19)
%% AASTeX v7.0.2 (aastex702.cls, 2025/05/09)
%% Compiled class files are vendored in this directory.
\documentclass{aastex702}

\shorttitle{UHI v5.2}
\shortauthors{Viridis}

%% Do not guess dates; leave submission dates unset.
\submitjournal{RNAAS}

\begin{document}

\title{A Probabilistic Unified Habitability Index Coupling Planetary Viability with JWST and ELT/PCS Observability across 6,112 Exoplanets}

\author[0009-0004-5620-5352,gname=Effy,sname=Viridis]{Effy Viridis}
\affiliation{Independent Researcher, El Salvador}
\email[show]{effywhitedva@gmail.com}
\correspondingauthor{Effy Viridis}

\begin{abstract}
Traditional habitability metrics such as the Earth Similarity Index (ESI) rank exoplanets by bulk physical analogy without accounting for atmospheric characterization feasibility. I present version 5.2-final of the Unified Habitability Index (UHI), a probabilistic, hurdle-based framework evaluated via 10,000 Monte Carlo draws per planet across 6,112 scored exoplanets from the NASA Exoplanet Archive. UHI combines seven physical, climatic, stellar, and instrumental sub-indices with Dirichlet weight propagation, synchronous M-dwarf climate parameterizations, a hybrid runaway greenhouse veto, and physical observability models for JWST and ELT/PCS. At a $5\sigma$ detection threshold, 515 planets remain viable ($P_{\rm alive} > 0$). Synthetic Solar System anchors at 10\,pc pass cleanly (Earth at the 99.30th (non-transiting) and 99.26th (transiting) percentiles; Venus and Mars $= 0$). Compared against ESI, UHI filters out 71.2\% of high-ESI ($\geq 0.8$) targets that are currently unobservable while elevating seven nearby non-transiting worlds into the Top 20 via reflected-light imaging.
\end{abstract}

%% ``Monte Carlo methods'' has no exact computational UAT match
%% (UAT 2238 is a laboratory reaction-model concept). Closest: Computational methods (1965).
\keywords{\uat{Astrobiology}{74} --- \uat{Computational methods}{1965} --- \uat{Direct imaging}{387} --- \uat{Exoplanets}{498} --- \uat{Habitable planets}{695} --- \uat{Transmission spectroscopy}{2133}}

\section{Introduction}

Metrics such as the Earth Similarity Index \citep[ESI;][]{SchulzeMakuch2011} and Habitable Zone limits \citep{Kopparapu2013} ignore characterization feasibility, so many top-ranked analogs lie beyond current or near-future observatories. I summarize UHI v5.2-final, an open Monte Carlo framework coupling intrinsic habitability with JWST (transiting) and ELT/PCS (non-transiting) observability.

\section{Method}

I score 6,112 planets from a 2026-11-18 snapshot of the NASA Exoplanet Archive \texttt{pscomppars} table \citep{Christiansen2025} with 10,000 Monte Carlo draws per planet (seed 42), with split-normal uncertainties and \citet{Chen2017} masses or radii where missing. UHI is a weighted geometric mean, ${\rm UHI} = \prod S_i^{w_i}$, of seven sub-indices with weights drawn from ${\rm Dirichlet}(2,2,2,2,1,1,2)$; any $S_i \leq 10^{-12}$ vetoes the planet (${\rm UHI} = 0$).

Atmospheric retention follows the bolometric Cosmic Shoreline \citep{Zahnle2017}, while all detection terms reside in the observability term: JWST transmission (TSM-scaled; \citealt{Kempton2018}) plus emission for transiting planets, and a photon-noise-limited ELT/PCS reflected-light model \citep{Kasper2021} otherwise. The climate term combines a 1-D energy balance model \citep{North1981} and $T_{\rm eff}$-dependent albedos \citep{Shields2013} with GCM-calibrated parameterizations for synchronous M-dwarf planets ($T_{\rm eff} < 3800$\,K): substellar clouds phased in by a flux ramp \citep{Yang2013}, an expanded inner edge \citep{Kopparapu2016}, and a CO$_2$ outer-edge floor \citep{Turbet2018}. Other planets face a runaway-greenhouse veto above 282\,W\,m$^{-2}$ absorbed flux \citep{Goldblatt2013,Goldblatt2014}. Mean molecular weight transitions log-sigmoidally across the radius valley \citep{Fulton2017}. Because vetoes make posteriors bimodal, I rank by ${\rm UHI}_{\rm rank} = P_{\rm alive} \times {\rm median}({\rm UHI} \mid {\rm UHI} > 0)$.

\section{Results}

At both $5\sigma$ (primary) and $3\sigma$ thresholds, 515 planets remain viable ($P_{\rm alive} > 0$; Spearman $\rho = 0.985$ between variants). Solar System twins at 10\,pc validate the pipeline: Earth ranks at the 99.30th (\#44, non-transiting) and 99.26th (\#46, transiting) percentiles with $P_{\rm alive} = 1.00$, while Venus and Mars score ${\rm UHI} = 0$.

At $5\sigma$ (Table~\ref{tab:top10}), temperate M-dwarf sub-Neptunes lead (TOI-4336 A b, 0.562; TOI-7166 b, 0.547), favored by large H$_2$-rich scale heights and the 2.5\,$R_\oplus$ upper radius limit $f_{R,d}$. At $3\sigma$, a proxy for deeper JWST campaigns, rocky worlds rise: TRAPPIST-1 d from \#5 (0.436) to \#1 (0.609) and TRAPPIST-1 e from \#17 (0.272) to \#9 (0.428). This matches a 51-parameter Sobol analysis \citep[SALib;][]{Herman2017} from v5 (not rerun; means over 16 outputs): the detection threshold (mean $S_T = 0.45$) and $f_{R,d}$ (0.15) dominate, the radius-valley and inner-edge transition parameters have $S_T < 0.01$, and the runaway threshold and cloud ramp were not sampled.

\begin{deluxetable}{cclcccccc}
\tablecaption{Top 10 UHI targets at $5\sigma$, with $3\sigma$ ranks, conditional posteriors, and ESI$_{\rm SM11}$ \citep{SchulzeMakuch2011}.\label{tab:top10}}
\tablewidth{0pt}
\tabletypesize{\scriptsize}
\tablehead{
\colhead{Rank ($5\sigma$)} & \colhead{Rank ($3\sigma$)} & \colhead{Planet} &
\colhead{UHI$_{\rm rank}$ ($5\sigma$)} & \colhead{$P_{\rm alive}$} &
\colhead{Cond.\ median [IQR]} & \colhead{$P$(top 10)} & \colhead{Tier} & \colhead{ESI$_{\rm SM11}$}
}
\startdata
1 & 2 & TOI-4336 A b & 0.562 & 0.99 & 0.567 [0.460--0.665] & 0.60 & T0 & 0.741 \\
2 & 3 & TOI-7166 b & 0.547 & 0.99 & 0.554 [0.456--0.647] & 0.56 & T1 & 0.816 \\
3 & 4 & GJ 273 b & 0.477 & 0.87 & 0.546 [0.424--0.642] & 0.47 & T2 & 0.847 \\
4 & 6 & K2-332 b & 0.449 & 0.91 & 0.494 [0.393--0.584] & 0.36 & T2 & 0.786 \\
5 & 1 & TRAPPIST-1 d & 0.436 & 1.00 & 0.436 [0.261--0.602] & 0.35 & T0 & 0.905 \\
6 & 8 & Wolf 1061 c & 0.434 & 0.82 & 0.529 [0.426--0.625] & 0.41 & T2 & 0.823 \\
7 & 11 & Proxima Cen d & 0.389 & 0.84 & 0.463 [0.356--0.566] & 0.28 & T2 & 0.876 \\
8 & 5 & TOI-2094 b & 0.377 & 0.99 & 0.382 [0.229--0.521] & 0.23 & T2 & 0.813 \\
9 & 17 & GJ 1061 d & 0.342 & 0.64 & 0.537 [0.429--0.611] & 0.33 & T2 & 0.884 \\
10 & 15 & K2-125 b & 0.333 & 0.84 & 0.396 [0.271--0.502] & 0.17 & T2 & 0.767 \\
\enddata
\end{deluxetable}

UHI diverges from ESI$_{\rm SM11}$ and the Habitable Worlds Catalog \citep{Mendez2024} (Spearman $\rho = 0.466$ and 0.444 among viable worlds). Of planets with ESI$_{\rm SM11} \geq 0.8$, 71.2\% (47/66) have median $P_{\rm det} < 0.001$ at $5\sigma$ (PHL ESI: 72.0\%, 18/25); none enter the top 29 (best: Kepler-1653 b, \#30), while 26 nearby non-transiting viable worlds reach median $P_{\rm det} \geq 0.5$ with ELT/PCS, seven in the Top 20 (GJ 273 b, Wolf 1061 c, Proxima Cen d and b, GJ 1061 d, Ross 128 b, GJ 433 d).

Sweeping the runaway threshold over 275--295\,W\,m$^{-2}$ keeps $P_{\rm alive}({\rm Venus}) = 0$, the Top 10 unchanged, and $\rho \geq 0.996$. A 3700--3750\,K $T_{\rm eff}$ veto boundary only swaps \#10 K2-125 b (3654\,K, 0.333) with \#11 Ross 128 b (0.323); I retain 3800\,K without post-hoc tuning.

Code (MIT), full methods, data, and figures (CC BY 4.0) are at \url{https://github.com/effywhitedva-ops/uhi-habitability-index} and Zenodo (\dataset[doi:10.5281/zenodo.24358996]{https://doi.org/10.5281/zenodo.24358996}); PHL data are not redistributed.

\begin{acknowledgments}
This research has made use of the NASA Exoplanet Archive, which is operated by the California Institute of Technology, under contract with the National Aeronautics and Space Administration under the Exoplanet Exploration Program. \textbf{AI use:} two custom agents built on Grok Bot (Shoggothsito and Aetherios) wrote and refactored code, ran simulations (local and Google Colab), produced tables and figures, audited methods, and assisted with literature search and editing. Gemini reviewed methods and drafts. I verified every step, checked all references against DOI records and NASA ADS, and take full responsibility for the content.
\end{acknowledgments}

\software{Python 3.13, NumPy, SciPy, pandas, SALib, Matplotlib, Astropy, PyArrow}

\begin{thebibliography}{99}

\bibitem[Chen \& Kipping(2017)]{Chen2017} Chen, J., \& Kipping, D. 2017, ApJ, 834, 17

\bibitem[Christiansen et al.(2025)]{Christiansen2025} Christiansen, J.~L., McElroy, D.~L., Harbut, M., et al. 2025, PSJ, 6, 186

\bibitem[Fulton et al.(2017)]{Fulton2017} Fulton, B.~J., Petigura, E.~A., Howard, A.~W., et al. 2017, AJ, 154, 109

\bibitem[Goldblatt et al.(2013)]{Goldblatt2013} Goldblatt, C., Robinson, T.~D., Zahnle, K.~J., \& Crisp, D. 2013, NatGeo, 6, 661

\bibitem[Goldblatt et al.(2014)]{Goldblatt2014} Goldblatt, C., Robinson, T.~D., Zahnle, K.~J., \& Crisp, D. 2014, NatGeo, 7, 931 (Corrigendum)

\bibitem[Herman \& Usher(2017)]{Herman2017} Herman, J., \& Usher, W. 2017, JOSS, 2, 97

\bibitem[Kasper et al.(2021)]{Kasper2021} Kasper, M., Cerpa Urra, N., Pathak, P., et al. 2021, Msngr, 182, 38

\bibitem[Kempton et al.(2018)]{Kempton2018} Kempton, E.~M.-R., Bean, J.~L., Louie, D.~R., et al. 2018, PASP, 130, 114401

\bibitem[Kopparapu et al.(2013)]{Kopparapu2013} Kopparapu, R.~K., Ramirez, R., Kasting, J.~F., et al. 2013, ApJ, 765, 131

\bibitem[Kopparapu et al.(2016)]{Kopparapu2016} Kopparapu, R.~K., Wolf, E.~T., Haqq-Misra, J., et al. 2016, ApJ, 819, 84

\bibitem[{PHL @ UPR Arecibo}(2024)]{Mendez2024} PHL @ UPR Arecibo 2024, The Habitable Worlds Catalog (file last updated 2024-03-12), \url{https://phl.upr.edu/hwc}

\bibitem[North et al.(1981)]{North1981} North, G.~R., Cahalan, R.~F., \& Coakley, J.~A. 1981, RGeo, 19, 91

\bibitem[Schulze-Makuch et al.(2011)]{SchulzeMakuch2011} Schulze-Makuch, D., M\'endez, A., Fair\'en, A.~G., et al. 2011, Astrobiology, 11, 1041

\bibitem[Shields et al.(2013)]{Shields2013} Shields, A.~L., Meadows, V.~S., Bitz, C.~M., et al. 2013, Astrobiology, 13, 715

\bibitem[Turbet et al.(2018)]{Turbet2018} Turbet, M., Bolmont, E., Leconte, J., et al. 2018, A\&A, 612, A86

\bibitem[Yang et al.(2013)]{Yang2013} Yang, J., Cowan, N.~B., \& Abbot, D.~S. 2013, ApJL, 771, L45

\bibitem[Zahnle \& Catling(2017)]{Zahnle2017} Zahnle, K.~J., \& Catling, D.~C. 2017, ApJ, 843, 122

\end{thebibliography}

\end{document}
